% !TEX program = lualatex
\documentclass[aspectratio=43,professionalfonts,handout]{beamer}
\usepackage{beamer_template}
\usepackage{enumerate}

\newcommand{\LectureNum}{11}
\newcommand{\LectureTitle}{複素フーリエ級数}
\newcommand{\CourseName}{応用数学}
\renewcommand{\TermName}{前期}

\begin{document}

\begin{frame}{第11回　複素フーリエ級数}
  \begin{block}{本回の内容}
  複素フーリエ級数
  \end{block}
  \vfill\hfill{\scriptsize \textcolor{gray}{第3章 フーリエ解析　§1}}
\end{frame}

\begin{frame}{複素フーリエ級数}
  \begin{block}{}
  {\large 複素フーリエ級数}
  \end{block}
  \vfill\hfill{\scriptsize \textcolor{gray}{p.\,86--89}}
\end{frame}

\begin{frame}{定義：複素フーリエ展開}
  \begin{block}{}
  \[
    f(t) = \sum_{n=-\infty}^{\infty} c_{n} e^{int}
  \]
  {\footnotesize \textcolor{gray}{オイラーの公式 $e^{int} = \cos(nt) + i \sin(nt)$ を利用}}
  \end{block}
  \vfill\hfill{\scriptsize \textcolor{gray}{p.\,86--89}}
\end{frame}

\begin{frame}{一般周期関数のフーリエ展開}
  \begin{block}{}
  \[
    f(t) = \sum_{n=-\infty}^{\infty} c_{n} e^{in\pi t/l}
  \]
  \end{block}
\end{frame}

\begin{frame}{例題6　（p.\,88）}
  \begin{exampleblock}{問題}
    次の周期 4 の関数 $f(x)$ の複素フーリエ級数を求めよ\\[2pt]
    $f(x) = x + 1  (-2 \leq x < 2),  f(x+4) = f(x)$\\[2pt]
  \end{exampleblock}
\end{frame}

\begin{frame}{例題6　解答}
  \begin{exampleblock}{解答}
  \begin{enumerate}[1.]
    \item[] $c_{n} = (1/4) \int_{-2}^{2} (x+1) e^{-in\pi x/2} dx$
    \item[] n \neq 0 のとき $: e^{in\pi} = (-1)^n$ を用いて計算
    \item[] $c_{n} = 2(-1)^n / (in\pi)$
    \item[] $n = 0$ のとき $: c_{0} = (1/4)\int_{-2}^{2}(x+1)dx = 1$
  \end{enumerate}
  \end{exampleblock}
  \begin{alertblock}{答}
    $f(x) = 1 + \sum_{n\neq0} 2(-1)^n/(in\pi) \cdot e^{in\pi x/2}$
  \end{alertblock}
\end{frame}

\begin{frame}{演習問題（p.\,86--89）}
  \begin{exampleblock}{自分で解いてみよう}
  \begin{description}
    \item[問6] 次の周期 2 の関数 $f(x), g(x)$ の複素フーリエ級数を求めよ
    $(1) f(x) = { 0  (-1 \leq x < 0),  1  (0 \leq x < 1) },  f(x+2) = f(x)$\\
    $(2) g(x) = { 0  (-1 \leq x < 0),  x  (0 \leq x < 1) },  g(x+2) = g(x)$\\
    \item[問7] 次の周期 4 の関数 $f(x), g(x)$ の複素フーリエ級数を求めよ
    $(1) f(x) = { 2  (-2 \leq x < 0),  0  (0 \leq x < 2) },  f(x+4) = f(x)$\\
    $(2) g(x) = { 0  (-2 \leq x < 0),  2-x  (0 \leq x < 2) },  g(x+4) = g(x)$\\
  \end{description}
  \end{exampleblock}
\end{frame}

\begin{frame}{練習問題1　（p.\,90）}
  \begin{block}{}
  \begin{description}
    \item[問1] 次の周期 2 の関数 $f(x)$ のフーリエ級数を求めよ
    $f(x) = { 1  (-1 \leq x < 0),  1-x  (0 \leq x < 1) },  f(x+2) = f(x)$\\
    \item[問2] グラフが次の図で示されている周期 4 の関数のフーリエ級数を求めよ
    (1) 三角波 $: f(x) = 2 - |x| (-2 \leq x < 2), f(x+4) = f(x)$ （頂点 $y=2,$ 周期 4 ）\\
    (2) のこぎり波 $: f(x) = x (-1 \leq x < 1), f(x+2) = f(x)$ （周期 2 ）\\
    \item[問3] 次の関数について、以下の問いに答えよ
    $f(x) = x^{2}  (-1 \leq x < 1),  f(x+2) = f(x)$\\
    (1) フーリエ級数を求めよ\\
    $(2) x=0$ および $x=1$ とおくことにより、次の公式を導け\\
    $1 - 1/2^{2} + 1/3^{2} - \ldots = \pi^{2}/12$\\
    $1 + 1/2^{2} + 1/3^{2} + \ldots = \pi^{2}/6$\\
    \item[問4] 次の関数 $f(x)$ の複素フーリエ級数を求めよ
    $f(x) = e^x  (-\pi/2 \leq x < \pi/2),  f(x+\pi) = f(x)$\\
  \end{description}
  \end{block}
\end{frame}

\end{document}
